\documentclass{article}

\usepackage{amsmath}

\title{Ch. 30 HW Solutions}
\author{Principles of Macroeconomics}
\date{\today}

\begin{document}

\maketitle

\section*{Quick Quiz}

\begin{enumerate}

\item (b) -- Fiat money is money that only has value because we've collectively decided it has value.

\item (c) -- Paying with a credit card isn't actually paying; it's promising to pay later. A promise to pay later isn't a reliable medium of exchange, unit of account, or store of value.

\item (a) -- The Fed was established by Congress.

\item (c) -- If it buys bonds, it pays the bond owners. When it pays the bond owners, it can introduce new money into the economy.

\item (c) -- If the bank doesn't loan out some of the deposits, then the money supply is unchanged because \$100 of currency is converted into \$100 of deposits. If the bank does loan out some of the deposits, the money supply increases because \$100 of currency is converted into more than \$100 of deposits. That's because fractional-reserve banking creates money.

\item (d) -- \[ \$120 \cdot \frac{1}{0.25} = \$480 \]

\item (c) -- The percent change in bank capital is $5 \cdot -10\% = -50\%$. A 50\% decrease in \$200 leaves \$100.

\end{enumerate}

\section*{Questions for Review}

\begin{enumerate}

\item Money is a medium of exchange and a unit of account.

\item Commodity money has intrinsic value. Fiat money does not. Fiat money only has value because we've decreed it has value. We use fiat money

\item Demand deposits are deposits that you can withdraw whenever you want. Typically, they're checking account deposits. They're included in the money stock because you can use a debit card to pay with deposits just as easily as you can use cash.

\item The Fed sets monetary policy. The members of the Fed who vote on monetary policy are split between Presidential appointees and regionally appointed officials.

\item It buys government bonds.

\item If they lend out some of their deposits, they can charge interest on it, and earn a profit. Reserves and the money supply are related as follows:
	\begin{gather*}
	S_M = C + D \\
	D = \text{res} \cdot MM = \text{res} \cdot \frac{1}{\text{res ratio}} \\
	S_M = C + \text{res} \cdot \frac{1}{\text{res ratio}} 
	\end{gather*}

\end{enumerate}

\section*{Problems and Applications}

\begin{enumerate}

\item

	\begin{enumerate}
	
	\item Money -- Pennies are currency. They're accepted as a medium of exchange. Prices are denominated in dollars and cents, so they're a unit of account. And they store value because you can save a penny today and spend it in the future. 
	
	\item No, Mexican pesos are not a medium of exchange in the U.S. because they're not regularly accepted as payment. Similarly they are not the unit of account because prices are quoted in dollars. That said, they are a store of value because Mexican pesos are assets that can always be exchanged for U.S. dollars.
	
	\item No, a Picasso painting is not the unit of account in the U.S. because prices aren't quoted in Picasso paintings. Similarly, it's not a medium of exchange because Picasso paintings aren't regularly used to buy things. It is a store of value because it will have value in the future.
	
	\item No, a plastic credit card is not a medium of exchange because it is not a payment method; it’s a method of deferring payment. A credit card payment is a store of value -- although not a reliable one -- because it is a promise to pay later. It’s not a unit of account because prices are denoted in today's dollars, not future dollar payments.
	
	\end{enumerate}

\end{enumerate}

\end{document}